% References cited by the book-level front and back matter. Exact DOI/URL
% identifiers merge with chapter entries in the single bibliography.
\begin{thebibliography}{99}
\bibitem{dnacomputingbook} \emph{DNA Computing: Molecular mechanisms, algorithms and
evidence}. Companion volume, prepublication review edition (October 2026).
\url{https://github.com/hermonai/dna-computing-book}.
\bibitem{raschkaoverview} S. Raschka. \emph{Build a Large Language Model (From Scratch)}:
author companion overview and chapter map.
\url{https://www.sebastianraschka.com/llms-from-scratch/}.
\bibitem{adleman1994} L. M. Adleman. Molecular computation of solutions to combinatorial
problems. \emph{Science} 266 (1994), 1021--1024.
\url{https://doi.org/10.1126/science.7973651}.
\bibitem{zhangwinfree2009} D. Y. Zhang and E. Winfree. Control of DNA strand displacement
kinetics using toehold exchange. \emph{Journal of the American Chemical Society} 131 (2009),
17303--17314. \url{https://doi.org/10.1021/ja906987s}.
\bibitem{soloveichik2010} D. Soloveichik, G. Seelig and E. Winfree. DNA as a universal
substrate for chemical kinetics. \emph{Proceedings of the National Academy of Sciences USA}
107 (2010), 5393--5398. \url{https://doi.org/10.1073/pnas.0909380107}.
\bibitem{qianwinfree2011} L. Qian and E. Winfree. Scaling up digital circuit computation
with DNA strand displacement cascades. \emph{Science} 332 (2011), 1196--1201.
\url{https://doi.org/10.1126/science.1200520}.
\bibitem{qianwinfreebruck2011} L. Qian, E. Winfree and J. Bruck. Neural network computation
with DNA strand displacement cascades. \emph{Nature} 475 (2011), 368--372.
\url{https://doi.org/10.1038/nature10262}.
\bibitem{cherryqian2018} K. M. Cherry and L. Qian. Scaling up molecular pattern recognition
with DNA-based winner-take-all neural networks. \emph{Nature} 559 (2018), 370--376.
\url{https://doi.org/10.1038/s41586-018-0289-6}.
\bibitem{cherryqian2025} K. M. Cherry and L. Qian. Supervised learning in DNA neural
networks. \emph{Nature} 645 (2025), 639--647.
\url{https://doi.org/10.1038/s41586-025-09479-w}.
\bibitem{songqian2025} T. Song and L. Qian. Heat-rechargeable computation in DNA logic
circuits and neural networks. \emph{Nature} 646 (2025), 315--322.
\url{https://doi.org/10.1038/s41586-025-09570-2}.
\bibitem{rothemund2006} P. W. K. Rothemund. Folding DNA to create nanoscale shapes and
patterns. \emph{Nature} 440 (2006), 297--302.
\url{https://doi.org/10.1038/nature04586}.
\bibitem{woods2019} D. Woods, D. Doty, C. Myhrvold, J. Hui, F. Zhou, P. Yin and E. Winfree.
Diverse and robust molecular algorithms using reprogrammable DNA self-assembly.
\emph{Nature} 567 (2019), 366--372. \url{https://doi.org/10.1038/s41586-019-1014-9}.
\bibitem{evans2024} C. G. Evans, J. O'Brien, E. Winfree and A. Murugan. Pattern recognition
in the nucleation kinetics of non-equilibrium self-assembly. \emph{Nature} 625 (2024),
500--507. \url{https://doi.org/10.1038/s41586-023-06890-z}.
\bibitem{church2012} G. M. Church, Y. Gao and S. Kosuri. Next-generation digital
information storage in DNA. \emph{Science} 337 (2012), 1628.
\url{https://doi.org/10.1126/science.1226355}.
\bibitem{goldman2013} N. Goldman, P. Bertone, S. Chen, C. Dessimoz, E. M. LeProust,
B. Sipos and E. Birney. Towards practical, high-capacity, low-maintenance information
storage in synthesized DNA. \emph{Nature} 494 (2013), 77--80.
\url{https://doi.org/10.1038/nature11875}.
\bibitem{erlich2017} Y. Erlich and D. Zielinski. DNA Fountain enables a robust and
efficient storage architecture. \emph{Science} 355 (2017), 950--954.
\url{https://doi.org/10.1126/science.aaj2038}.
\bibitem{evo2} G. Brixi et al. \emph{Genome modelling and design across all domains of life
with Evo 2}. Nature 652 (2026), 1349--1361; published 4 March 2026.
\url{https://doi.org/10.1038/s41586-026-10176-5}.
\bibitem{dnabert2021} Y. Ji, Z. Zhou, H. Liu and R. V. Davuluri.
\emph{DNABERT: pre-trained Bidirectional Encoder Representations from Transformers
model for DNA-language in genome}. Bioinformatics 37 (2021), 2112--2120.
\url{https://doi.org/10.1093/bioinformatics/btab083}.
\bibitem{dnabert2} Z. Zhou et al.
\emph{DNABERT-2: Efficient Foundation Model and Benchmark For Multi-Species Genome}.
arXiv:2306.15006 (2023; version 2, 18 March 2024; ICLR 2024).
\url{https://arxiv.org/abs/2306.15006}.
\bibitem{enformer2021} \v{Z}. Avsec et al.
Effective gene expression prediction from sequence by integrating long-range
interactions. \emph{Nature Methods} 18 (2021), 1196--1203.
\url{https://doi.org/10.1038/s41592-021-01252-x}.
\bibitem{nt2025} H. Dalla-Torre et al.
\emph{Nucleotide Transformer: building and evaluating robust foundation models
for human genomics}. Nature Methods 22 (2025), 287--297.
\url{https://doi.org/10.1038/s41592-024-02523-z}.
\bibitem{hyenadna} E. Nguyen et al.
\emph{HyenaDNA: Long-Range Genomic Sequence Modeling at Single Nucleotide Resolution}.
arXiv:2306.15794 (2023; version 2, 14 November 2023).
\url{https://arxiv.org/abs/2306.15794}.
\bibitem{caduceus} Y. Schiff et al.
\emph{Caduceus: Bi-Directional Equivariant Long-Range DNA Sequence Modeling}.
arXiv:2403.03234 (2024; version 2, 5 June 2024; ICML 2024).
\url{https://arxiv.org/abs/2403.03234}.
\bibitem{evo2024} E. Nguyen, M. Poli et al. Sequence modeling and design from molecular
to genome scale with Evo. \emph{Science} 386 (2024), eado9336.
\url{https://doi.org/10.1126/science.ado9336}.
\bibitem{alphagenome2026} \v{Z}. Avsec, N. Latysheva, J. Cheng et al. Advancing
regulatory variant effect prediction with AlphaGenome. \emph{Nature} 649 (2026),
1206--1218. \url{https://doi.org/10.1038/s41586-025-10014-0}.
\bibitem{evod48:benchmark} H. Feng et al. \emph{Benchmarking DNA foundation
models for genomic and genetic tasks}. Nature Communications (2025).
\url{https://doi.org/10.1038/s41467-025-65823-8}.
\bibitem{evod51:levels} M. R. Morris et al. \emph{Levels of AGI:
Operationalizing Progress on the Path to AGI}. 2023; arXiv v2.
\url{https://arxiv.org/abs/2311.02462}.
\bibitem{albertscontrol} B. Alberts et al. An overview of gene control. In
\emph{Molecular Biology of the Cell}, 4th edition. Garland Science (2002).
\url{https://www.ncbi.nlm.nih.gov/books/NBK26885/}.
\bibitem{ucscformats} UCSC Genome Browser. \emph{Data file formats} (BED definition:
first base numbered 0; \texttt{chromEnd} not included). Accessed 18 October 2026.
\url{https://genome.ucsc.edu/FAQ/FAQformat.html}.
\bibitem{vcfspec} SAMtools/hts-specs. \emph{The Variant Call Format specification,
VCFv4.3} (POS: the 1st base has position 1). Accessed 18 October 2026.
\url{https://samtools.github.io/hts-specs/VCFv4.3.pdf}.
\bibitem{gff3spec} Sequence Ontology. \emph{Generic Feature Format version 3
(GFF3)} (one-based start/end, start $\le$ end). Accessed 18 October 2026.
\url{https://github.com/The-Sequence-Ontology/Specifications/blob/master/gff3.md}.
\end{thebibliography}
